% UR3e XR Simulation — Mathematical Reference
% pdflatex math_reference.tex  (run twice for TOC)
\documentclass[11pt,a4paper]{article}

\usepackage{amsmath,amssymb}
\usepackage{geometry}
\usepackage{hyperref}
\usepackage{booktabs}
\usepackage{parskip}
\geometry{margin=2.5cm}

\title{UR3e XR Simulation\\[4pt]\large Mathematical Reference}
\author{}
\date{}

\begin{document}
\maketitle
\tableofcontents
\newpage

\section{Homogeneous Transforms (SE(3))}

Poses are $4\times4$ rigid-body matrices:
\[
T = \begin{bmatrix} R & \mathbf{t} \\ \mathbf{0}^\top & 1 \end{bmatrix},
\quad R \in SO(3),\; \mathbf{t} \in \mathbb{R}^3
\]
\[
T_{A\to C} = T_{A\to B}\,T_{B\to C}, \qquad
T^{-1} = \begin{bmatrix} R^\top & -R^\top\mathbf{t} \\ \mathbf{0}^\top & 1 \end{bmatrix}
\]

\section{Geometric Jacobian}

Maps joint velocities to EEF twist $[\dot{\mathbf{p}};\,\boldsymbol{\omega}]$:
\[
J = \begin{bmatrix}
\mathbf{z}_1 \times \mathbf{r}_1 & \cdots & \mathbf{z}_6 \times \mathbf{r}_6 \\[3pt]
\mathbf{z}_1 & \cdots & \mathbf{z}_6
\end{bmatrix}, \qquad
\mathbf{r}_i = \mathbf{p}_{\text{EEF}} - \mathbf{p}_i
\]
$\mathbf{z}_i$ is joint $i$'s axis in world space; $\mathbf{p}_i$ is its origin.

\section{RMRC — Resolved Motion Rate Control}

\textit{Whitney, 1969.} Proportional controller on Cartesian error:
\[
\mathbf{v}_{\text{des}} =
\begin{bmatrix} K_p\,\Delta\mathbf{p} \\ K_o\,w_\text{ori}\,\theta_e\,\hat{\mathbf{n}}_e \end{bmatrix}
\in \mathbb{R}^6, \qquad \dot{\mathbf{q}} = J^+\,\mathbf{v}_{\text{des}}
\]
$\theta_e, \hat{\mathbf{n}}_e$ are the angle and axis of $q_e = q_{\text{target}}\,q_{\text{EEF}}^{-1}$, wrapped to $[-\pi,\pi]$.

\section{Damped Least Squares (DLS)}

\textit{Nakamura \& Hanafusa, 1986.} Regularises the pseudoinverse near singularities:
\[
J^+_{\lambda} = \bigl(J^\top J + \lambda^2 I\bigr)^{-1} J^\top, \qquad
\dot{\mathbf{q}}_{\text{DLS}} = J^+_\lambda\,\mathbf{v}_{\text{des}}
\]
The $6\times6$ system is solved by Gaussian elimination with partial pivoting.

\section{Gradient Descent Fallback}

Deep in a singularity the DLS inverse is unreliable. A step size $\alpha$ replaces it:
\[
\dot{\mathbf{q}}_{\text{GD}} = \alpha\,J^\top \mathbf{v}_{\text{des}}
\]

\section{Manipulability \& DLS/GD Blend}

\textit{Yoshikawa, 1985.}
\[
w = |\det J|, \qquad
\beta = \frac{w}{w + w_\varepsilon}, \qquad
\dot{\mathbf{q}} = \beta\,\dot{\mathbf{q}}_{\text{DLS}} + (1-\beta)\,\dot{\mathbf{q}}_{\text{GD}}
\]
$\beta \to 1$ when well-conditioned (DLS dominates); $\beta \to 0$ near singular (GD takes over).

\section{Joint Velocity Limiting \& Euler Integration}

Scale the whole vector proportionally so the fastest joint stays within $\dot{q}_{\max}$:
\[
s = \min\!\left(1,\;\frac{\dot{q}_{\max}}{\max_j |\dot{q}_j|}\right), \qquad
\dot{\mathbf{q}} \leftarrow s\,\dot{\mathbf{q}}
\]
\[
\mathbf{q}_{k+1} = \mathbf{q}_k + \dot{\mathbf{q}}\,\Delta t
\]

\section{Singularity Detection}

\textit{Villalobos et al., 2022.} Closed-form UR3e Jacobian determinant:
\[
\det J = \sin q_3 \cdot \sin q_5 \cdot
\underbrace{(a_2 \cos q_2 + a_3 \cos(q_2{+}q_3) + d_5 \sin(q_2{+}q_3{+}q_4))}_{\text{shoulder term}}
\]

\begin{center}
\begin{tabular}{lll}
\toprule
\textbf{Type} & \textbf{Condition} & \textbf{Meaning} \\
\midrule
Elbow    & $\sin q_3 \approx 0$ & Arm fully extended or folded \\
Wrist    & $\sin q_5 \approx 0$ & J4/J6 axes parallel \\
Shoulder & shoulder term $\approx 0$ & Wrist centre on base axis \\
\bottomrule
\end{tabular}
\end{center}

UR3e DH parameters: $a_2 = 0.244\,\text{m}$, $a_3 = 0.213\,\text{m}$, $d_5 = 0.085\,\text{m}$.

\section{Null-Space Self-Collision Avoidance}

\textit{Khatib, 1986.} When $n > m$ the general IK solution is:
\[
\dot{\mathbf{q}} = J^+\dot{\mathbf{r}} + \underbrace{(I - J^+J)}_{P}\,\dot{\mathbf{q}}_0
\]
$P$ is the null-space projector ($P^2 = P$). Any $\dot{\mathbf{q}}_0$ projected through $P$
pursues a secondary objective (e.g.\ collision avoidance) without disturbing the task.

\paragraph{LQO derivation.}
Minimise $H = \tfrac{1}{2}(\dot{\mathbf{q}} - \dot{\mathbf{q}}_0)^\top W\,(\dot{\mathbf{q}} - \dot{\mathbf{q}}_0)$
subject to $J\dot{\mathbf{q}} = \dot{\mathbf{r}}$. Via Lagrange multipliers:
\[
J^+_W = W^{-1}J^\top(J W^{-1}J^\top)^{-1}, \qquad
\dot{\mathbf{q}}^* = J^+_W\,\dot{\mathbf{r}} + (I - J^+_W J)\,\dot{\mathbf{q}}_0
\]
$W = I$ recovers the Moore–Penrose pseudoinverse.

\section{Pinhole Camera Projection}

Forward and inverse (back-projection):
\[
u = f_x\frac{X}{Z} + c_x, \quad v = f_y\frac{Y}{Z} + c_y
\qquad\Longleftrightarrow\qquad
X = \frac{u-c_x}{f_x}Z, \quad Y = \frac{v-c_y}{f_y}Z
\]
\[
\theta_H = 2\arctan\!\left(\frac{W}{2f_x}\right)
\]
Depth is the median over a $9\times9$ pixel patch (R16 texture: normalised $[0,1]
\times 65535$ to get mm, divided by 1000 for metres).

\section{Madgwick 6-DOF IMU Filter (Y-up)}

\textit{Madgwick et al., 2011.} State: unit quaternion $\mathbf{q} = [q_0,q_1,q_2,q_3]^\top$.

Objective function with gravity reference $\hat{\mathbf{g}} = [0,1,0]$:
\begin{align*}
f_1 &= 2(q_1 q_2 - q_0 q_3) - a_x \\
f_2 &= 1 - 2(q_1^2 + q_3^2) - a_y \\
f_3 &= 2(q_2 q_3 + q_0 q_1) - a_z
\end{align*}

Gradient $\hat{\mathbf{s}} = \nabla_{\mathbf{q}}\|\mathbf{f}\|^2 / \|\cdot\|$:
\[
s_0 = -2q_3 f_1 + 2q_1 f_3, \quad
s_1 = 2q_2 f_1 - 4q_1 f_2 + 2q_0 f_3
\]
\[
s_2 = 2q_1 f_1 + 2q_3 f_3, \quad
s_3 = -2q_0 f_1 - 4q_3 f_2 + 2q_2 f_3
\]

Integration:
\[
\dot{\mathbf{q}} = \tfrac{1}{2}\,\mathbf{q}\otimes[0,g_x,g_y,g_z] - \beta\hat{\mathbf{s}}, \qquad
\mathbf{q}_{k+1} = \frac{\mathbf{q}_k + \dot{\mathbf{q}}\,\Delta t}{\|\mathbf{q}_k + \dot{\mathbf{q}}\,\Delta t\|}
\]

\section{Coordinate Frame Transforms}

ROS optical $\to$ Unity (Y-flip): $\mathbf{p}_U = (x,-y,z)$

Unity world $\to$ ROS world:
\[
(x_R,\,y_R,\,z_R) = (z_U,\,-x_U,\,y_U), \qquad
q^R = (-q^U_z,\,q^U_x,\,-q^U_y,\,q^U_w)
\]

Camera $\to$ world: $\mathbf{p}_W = R_\text{cam}\mathbf{p}_c + \mathbf{t}_\text{cam}$

\section{Hand-Eye Transform}

Camera pose in EEF frame, read from the URDF hierarchy in simulation:
\[
\mathbf{p}_\text{local} = R_\text{EEF}^{-1}(\mathbf{p}_\text{cam} - \mathbf{p}_\text{EEF}), \qquad
R_\text{local} = R_\text{EEF}^{-1}R_\text{cam}
\]
On the real robot this is the $AX = XB$ calibration problem:
\[
T_{\text{EEF}_i\to\text{EEF}_j}\,X = X\,T_{\text{cam}_i\to\text{cam}_j}
\]

\section{AprilTag Position Fusion}

Tag at known world position $\mathbf{p}_\text{tag}^W$, detected at camera-local $\mathbf{p}_\text{tag}^c$:
\[
\mathbf{p}_\text{cam}^W = \mathbf{p}_\text{tag}^W - R_\text{cam}^W\,\mathbf{p}_\text{tag}^c
\]
Fuses tag-derived position with Madgwick orientation to give a full 6-DOF world pose.

\section{Arc Trajectory}

The target sweeps the outer ring of the work volume rather than cutting straight
across, keeping clear of the wrist-singularity zone at the centre.
In cylindrical coordinates $(r,\theta,y)$ around the robot base:
\[
\theta(t) = \theta_0 + (\theta_1-\theta_0)\,t, \quad
r = r_\text{arc}, \quad
y(t) = y_0 + (y_1-y_0)\,t, \quad t\in[0,1]
\]
\[
x = r_\text{arc}\sin\theta(t), \quad z = r_\text{arc}\cos\theta(t)
\]
Work volume: $r \in [r_{\min}, r_{\max}]$, $y \in [y_\text{table},\; y_\text{table}+h]$.


\section{Compute Shaders}

\subsection{Thread dispatch — 16$\times$16$\times$1}

The kernel is declared with a $16\times16\times1$ thread group:
\begin{verbatim}
[numthreads(16, 16, 1)]
\end{verbatim}

Each thread processes exactly one depth pixel $(u,v)$ identified by the built-in
\texttt{SV\_DispatchThreadID} $\mathbf{id} = (id_x, id_y, 0)$.
The CPU dispatches:
\[
\left\lceil \frac{W}{16} \right\rceil \times \left\lceil \frac{H}{16} \right\rceil \times 1
\quad \text{thread groups}
\]

giving at most $16\times16$ threads per group and a total thread count of
$16\lceil W/16\rceil \times 16\lceil H/16\rceil$ — always $\geq W{\times}H$.
Threads where $id_x \geq W$ or $id_y \geq H$ are culled by a
bounds guard and write nothing.

The flat output buffer index is:
\[
\text{index} = id_y \cdot W + id_x
\]

\subsection{Why $16\times16$ and not $8\times8$?}

The choice is driven by three hardware constraints.

\paragraph{1.\ Warp / wavefront alignment.}
NVIDIA GPUs execute threads in \emph{warps} of 32; AMD in \emph{wavefronts} of 64.
A thread group is partitioned into warps in row-major order.

\begin{center}
\renewcommand{\arraystretch}{1.3}
\begin{tabular}{lcc}
\toprule
\textbf{Config} & \textbf{Threads/group} & \textbf{Warps/group (NVIDIA)} \\
\midrule
$8\times8\times1$   & 64  & 2 \\
$16\times16\times1$ & 256 & 8 \\
\bottomrule
\end{tabular}
\end{center}

With only 2 warps per group the warp scheduler has little room to hide memory
latency (texture fetch $\sim$100--400 cycles).  8 warps keeps the execution
pipelines fed while earlier warps are stalled on reads.

\paragraph{2.\ Memory bus utilisation (coalescing).}
A row of 16 consecutive threads reads 16 adjacent depth pixels.
For a 16-bit (2-byte) R16 texture this is a single 32-byte transaction, exactly
filling one 32-byte memory sub-transaction on current hardware:
\[
16\;\text{threads} \times 2\;\text{bytes} = 32\;\text{bytes} = 1\;\text{transaction}
\]
An $8\times8$ group issues rows of 8 threads:
\[
8\;\text{threads} \times 2\;\text{bytes} = 16\;\text{bytes} \;\Rightarrow\; \text{half-utilised bus}
\]
Halving the row width doubles the number of memory transactions for the same pixel count.

\paragraph{3.\ Dispatch overhead at D455 resolution.}
For the D455 depth stream ($1280\times720$):

\[
16\times16:\quad \left\lceil\frac{1280}{16}\right\rceil \times \left\lceil\frac{720}{16}\right\rceil
= 80 \times 45 = 3\,600 \text{ groups}
\]
\[
8\times8:\quad \left\lceil\frac{1280}{8}\right\rceil \times \left\lceil\frac{720}{8}\right\rceil
= 160 \times 90 = 14\,400 \text{ groups}
\]

$14\,400$ groups incurs $4\times$ the GPU command-processor overhead with no
reduction in total thread count ($921\,600$ either way).

\medskip\noindent
\textbf{Summary:} $16\times16\times1$ maximises warp-level latency hiding,
fully coalesces 16-bit texture reads into single 32-byte transactions, and
reduces dispatch overhead by $4\times$ compared to $8\times8\times1$ — at the cost
of zero additional shared memory or register pressure for this kernel.

\subsection{GPU back-projection}

Rearranging the forward model $u = f_x(X/Z)+c_x$,\; $v = f_y(Y/Z)+c_y$:
\[
X = Z\bigl(id_x \cdot \tfrac{1}{f_x} - \tfrac{c_x}{f_x}\bigr), \quad
Y = Z\bigl(id_y \cdot \tfrac{1}{f_y} - \tfrac{c_y}{f_y}\bigr), \quad
Z = d_\text{raw}\times 65535 \times 0.001
\]
Pre-computed reciprocals (\texttt{invFx}, \texttt{invFy}, \texttt{cx\_over\_fx},
\texttt{cy\_over\_fy}) avoid per-thread division. Y is negated to flip ROS optical
(Y-down) to Unity (Y-up).

Colour pixel coordinates scale with image resolution ratio:
\[
u_c = \lfloor id_x \cdot W_c/W \rfloor, \qquad v_c = \lfloor id_y \cdot H_c/H \rfloor
\]

\subsection{Quad expansion (\texttt{PointCloudVertexColor.shader})}

A raw depth pixel maps to a single screen pixel — at short range the cloud looks
sparse and holey. Each point is instead expanded into a tiny screen-facing quad
(4 vertices, 2 triangles) in the CPU index buffer (\texttt{MeshTopology.Triangles}).

The vertex shader reads \texttt{SV\_VertexID}:
\[
\text{pointIdx} = \lfloor id/4 \rfloor, \qquad \text{cornerIdx} = id \bmod 4
\]

Half-size scales with depth so quads stay proportional to sensor pixel spacing:
\[
hs = s_q \cdot Z
\]

Corner offsets encode X/Y sign bits from \texttt{cornerIdx}:

\begin{center}
\begin{tabular}{ccc}
\toprule
\textbf{Corner} & \textbf{Bit pattern} & \textbf{Offset} \\
\midrule
0 (TL) & \texttt{00} & $(-hs,+hs)$ \\
1 (BL) & \texttt{01} & $(-hs,-hs)$ \\
2 (TR) & \texttt{10} & $(+hs,+hs)$ \\
3 (BR) & \texttt{11} & $(+hs,-hs)$ \\
\bottomrule
\end{tabular}
\end{center}

Each corner is projected to clip space individually so perspective foreshortening is exact.

Default $s_q = 0.003$: the D455 at 640$\times$480 has $\approx 0.134°$ per pixel, giving a
physical neighbour gap of $Z\tan(0.134°) \approx 0.00234\,Z$. Using $0.003Z$ adds
$\mathord{\sim}30\%$ overlap so adjacent quads always close without visible holes at any depth.

No geometry shader is used because URP's depth and shadow passes have no geometry stage;
a \texttt{MeshTopology.Points} + geometry shader combination causes a Vulkan topology
mismatch and the draw is rejected.

\subsection{TSDF integration}

Kernels \texttt{Integrate} and \texttt{Clear} use $8^3$ thread groups:
$\lceil N/8\rceil^3$ groups dispatched.

Voxel centre in world space:
\[
\mathbf{p}_W = \mathbf{c}_\text{vol} + \left([i,j,k]^\top - \tfrac{N}{2} + \tfrac{1}{2}\right)v_s, \qquad v_s = L/N
\]
Project to depth image: $\mathbf{p}_c = T_{W\to c}\,[\mathbf{p}_W;1]$,\; $u = f_x p_{cx}/p_{cz}+c_x$.

Signed distance and truncation:
\[
\text{SDF} = z_m - p_{cz}, \qquad
\text{TSDF} = \text{clamp}(\text{SDF}/d_t,\,-1,\,1)
\]
Voxels with $\text{SDF} < -d_t$ are skipped.

KinFu weighted running average:
\[
\tilde{d}_\text{new} = \frac{\tilde{d}_\text{old}\,w_\text{old} + \text{TSDF}}{w_\text{old}+1}, \qquad
w_\text{new} = \min(w_\text{old}+1,\;w_{\max})
\]

Surface extraction (\texttt{ExtractPoints}, $8^3$ groups): emit voxel $(i,j,k)$ if
\[
\text{sign}(\tilde{d}_{i,j,k}) \neq \text{sign}(\tilde{d}_{i+1,j,k})
\;\lor\; \tilde{d}_{i,j+1,k}
\;\lor\; \tilde{d}_{i,j,k+1}
\]
Slot allocated atomically: $\texttt{slot} = \texttt{InterlockedAdd}(\texttt{Counter}[0], 1)$.

Confidence colour:
\[
B = (1-|\tilde{d}|)\cdot\text{sat}(w/20), \qquad
\text{RGBA} = (0.7B,\;0.85B,\;0.7B+0.3\,\text{sat}(w/20),\;1)
\]

\end{document}

